\documentclass[aspectratio=169]{beamer}
\usetheme{default}
\usepackage{booktabs}
\definecolor{editblue}{RGB}{38,117,168}
\setbeamercolor{structure}{fg=editblue}
\title{How much edit experience should the model receive?}
\subtitle{Sequence-edit JEPA validity audit and bounded pilot}
\date{17 July 2026}
\begin{document}
\frame{\titlepage}

\begin{frame}{The task is synthetic typo repair}
  \begin{itemize}
    \item Official iGSM supplies a problem and its clean, multi-step worked solution.
    \item We damage that solution with token insertions, deletions, and replacements.
    \item The expert path is the exact reverse of those corruptions, so it is oracle denoising rather than naturally observed editing.
    \item More independent damaged solutions add breadth; more alternative edits per state add local action coverage.
  \end{itemize}
  \vfill
  \textbf{Why it matters:} we must establish healthy primitive dynamics before asking whether hierarchy helps.
\end{frame}

\begin{frame}{The old hierarchy round is invalid}
  \begin{itemize}
    \item All five jobs failed before their first optimizer step, so they provide no architectural evidence.
    \item The immediate symptom was a 276-token ``sentence'' exceeding a 128-position table.
    \item The cause was deeper: official steps often end in fused tokens such as \texttt{6.}, while the adapter split only on standalone \texttt{.}.
    \item The corrected adapter preserves official nested steps and exactly reverses every audited corruption.
  \end{itemize}
  \vfill
  \textbf{What to notice:} this is an invalid result, not evidence against hierarchy.
\end{frame}

\begin{frame}{The corrected data contract passes its first gate}
  \centering
  \begin{tabular}{lr}
    \toprule
    Audit check & Observed result \\
    \midrule
    Exact terminal recovery & 128/128 \\
    Multi-step solution collapse & 0/128 \\
    Official solution length & median 166.5 tokens \\
    Official steps & median 9 \\
    Maximum step length & p95 43 tokens \\
    Repairs per trajectory & mean 11.14 \\
    Repair/minimum-distance ratio & mean 1.029 \\
    \bottomrule
  \end{tabular}
  \vfill
  \textbf{Why it matters:} the model now sees the intended step sequence, but these are data-validity observations, not trained-model results.
\end{frame}

\begin{frame}{The next pilot varies one data axis at a time}
  \begin{itemize}
    \item The K=0 data curve uses 512, 2,000, and 6,000 unique trajectories for three passes.
    \item The counterfactual curve uses K=0, 1, 4, and 8 exact alternative outcomes at the common 2,000-trajectory anchor.
    \item Alternatives use observed current-buffer tokens and carry no preference, defect, or remaining-distance label.
    \item A 128-example H4 cell is only a full-shape process smoke; architecture comparisons remain gated.
    \item We retain the smallest K near the best result only if it improves held-out dynamics without factual-error or rank regression.
  \end{itemize}
  \vfill
  \textbf{Next decision:} spend later compute on unique diversity, local alternatives, or neither.
\end{frame}
\end{document}
